\documentclass[aspectratio=169, 11pt]{beamer}

% --- Packages & Setup ---
\usepackage{amsmath}
\usepackage{amssymb}
\usepackage[utf8]{inputenc}
\usepackage[T1]{fontenc}
\usepackage{xcolor}
\usepackage{booktabs} 
\usepackage{graphicx}
\usepackage{tikz} % Required for the corner box
\usetikzlibrary{calc}
\usetikzlibrary{positioning}
% Try loading Open Sans; fallback if missing
\IfFileExists{opensans.sty}{\usepackage[default,scale=0.95]{opensans}}{}

% --- Branding Colors ---
\definecolor{amritaMaroon}{HTML}{A4123F} 

% --- Theme Setup ---
\usetheme{Madrid}
\usecolortheme[named=amritaMaroon]{structure}
\setbeamertemplate{navigation symbols}{}
\setbeamertemplate{footline}[frame number]

\begin{document}

% --- CUSTOM TITLE SLIDE (Research) ---
\begin{frame}
    % Corner Box Indicator
    \begin{tikzpicture}[remember picture,overlay]
        \node[fill=amritaMaroon, text=white, font=\bfseries\footnotesize, anchor=north east, rounded corners=2pt, inner sep=5pt] 
        at ($(current page.north east) + (-0.2cm,-0.2cm)$) 
        {Type 1: Research Project};
    \end{tikzpicture}

    \centering
    \includegraphics[width=3.5cm]{logo.png} 
    \vspace{0.3cm}
    
    {\Large \textbf{\textcolor{amritaMaroon}{A Hierarchical Approach for Efficient Virtual Screening Using 2D Filtering and 3D Structural Similarity}}} \\
    \vspace{0.2cm}
    {\normalsize \textbf{Team Number: [11]}} \\ 
    \vspace{0.5cm}

    % Student Table
    \begin{table}
        \centering
        \small 
        \renewcommand{\arraystretch}{1.2}
        \begin{tabular}{|c|c|l|}
            \hline
            \textbf{S.No} & \textbf{Register Number} & \textbf{Student Name} \\
            \hline
            1 & CB.SC.U4CSE23206 & Allampati Chenchu Sanhitha Reddy \\
            \hline
            2 & CB.SC.U4CSE23245 & Shaik Abdus Sattar \\
            \hline
            3 & CB.SC.U4CSE23616 & Geethika G Nair \\
            \hline
            4 & CB.SC.U4CSE23633 & Meera J Vaishnav \\
            \hline
        \end{tabular}
    \end{table}
    \vspace{0.3cm}
    
    Guide: \textbf{Dr. T Ramraj} \\
    {\footnotesize Assistant Professor(SI.Gd), Department of CSE}
    
    \vspace{0.5cm}
    {\footnotesize \today}
\end{frame}

% --- SLIDE 1:PANEL COMMENTS  ---
\begin{frame}{1. Panel Comments}
\small
\begin{itemize}

\item \textbf{Computational Feasibility:} Direct 3D graph kernel methods are computationally expensive ($O(n^6)$), whereas our 2D topological screening runs at $O(T \cdot n^3)$ per pair --- competitive with Random Walk and Optimal Assignment kernels --- and a rank-one algebraic factorization (Lemma~1) eliminates one $O(n^3)$ sub-step from each gradient evaluation, yielding measured speeds of ${\sim}6$\,ms/pair on drug-like molecules.

\item \textbf{Novelty:} We propose a hybrid two staged (2D to 3D) pipeline where efficient 2D similarity screening identifies promising candidates, and only those are further analyzed using computationally intensive 3D kernels.

\item \textbf{Output \& Validation:} The model predicts molecular binding activity (active/inactive) and is evaluated using standard classification metrics.

\item \textbf{2D from 3D Data:} Although the dataset contains 3D molecular structures, we initially extract only 2D topological features such as atom connectivity and bond types, ignoring geometric properties.

\item \textbf{Why 2D to 3D:} The 2D stage efficiently reduces the search space by selecting highly similar molecules (e.g., similarity 0.9), after which full 3D structural features are incorporated for precise similarity computation.

\end{itemize}
\end{frame}

% --- SLIDE 2: PROBLEM STATEMENT ---
\section{Problem Statement}
\begin{frame}{2. Problem Statement}
\framesubtitle{Efficient Virtual Screening via Hierarchical 2D–3D Similarity Modeling}

\vspace{1em}

\begin{itemize}
\setlength\itemsep{1em}

\item \textbf{Domain \& Objective:} Virtual screening in drug discovery aims to identify molecules that can effectively bind to a target protein.

\item \textbf{Key Challenge:} Existing 3D similarity methods (e.g., graph kernels) are accurate but computationally expensive, requiring evaluation of all candidates at high complexity.

\item \textbf{Research Gap:} Current approaches lack an efficient mechanism to reduce the search space before applying costly 3D computations.

\item \textbf{Proposed Direction:} A hierarchical pipeline that performs fast 2D topological filtering followed by selective 3D kernel-based analysis.

\end{itemize}

\vspace{0.5em}

\centering
\textit{\small Goal: Achieve efficient and scalable molecular similarity prediction without compromising structural accuracy.}

\end{frame}

\section{Literature Survey}

% --- SLIDE 3 : Literature Survey---
\begin{frame}{3. Literature Survey}
\framesubtitle{Survey Scope}
\begin{itemize}

    \item We conducted a comprehensive survey of both foundational and recent research works in the domain of molecular property prediction and virtual screening.

    \item The survey primarily focuses on three major methodological directions:
    \begin{itemize}
        \item \textbf{Graph Kernel Methods:} Traditional approaches that compute similarity between molecular graphs using handcrafted structural features and kernel functions.
        
        \item \textbf{Graph Neural Networks (GNNs):} Deep learning-based approaches that automatically learn representations from molecular graphs, capturing complex patterns in data.
        
        \item \textbf{2D and 3D Molecular Representations:} Techniques that model molecules using:
        \begin{itemize}
            \item 2D topological structures (atoms and bonds)
            \item 3D geometric features (bond length, bond angle, torsion angle)
        \end{itemize}
    \end{itemize}

    \item The selected literature includes:
    \begin{itemize}
        \item \textbf{Classical approaches} such as Weisfeiler–Lehman graph kernels and pharmacophore kernels
        \item \textbf{Modern deep learning methods} such as GCNs and multimodal GNN architectures
        \item \textbf{Recent advancements} incorporating 3D geometric information and hybrid representations
    \end{itemize}
\end{itemize}
\end{frame}

%--------------------------------------------------
% --- TABLE SLIDE 1 ---
\begin{frame}{3.1 Survey: Key Papers with Research Gaps}
\scriptsize
\setlength{\tabcolsep}{2pt}

\begin{tabular}{|p{2.8cm}|p{2.3cm}|p{0.9cm}|p{3.2cm}|p{3.8cm}|}
\hline
\textbf{Paper Title} & \textbf{Authors} & \textbf{Year} & \textbf{Focus} & \textbf{Research Gap} \\ \hline

A multimodal geometric graph neural network approach for molecular property prediction \cite{jiang2026qegnn}
& Zhixin Jiang et al. 
& 2026 
& Multimodal fusion of fingerprints, 2D + 3D features 
& Computationally expensive; requires large training data \\ \hline

Self-conformation-aware framework \cite{qiao2025self}
& Jianbo Qiao et al. 
& 2025 
& Conformation-aware 2D representation learning 
& Does not fully utilize complete 3D geometry \\ \hline

Multilevel Fusion GNN \cite{liu2025multilevel}
& XiaYu Liu et al. 
& 2025 
& Hierarchical modeling of 2D graph structures 
& No explicit modeling of 3D spatial features \\ \hline

Weisfeiler–Lehman Kernel \cite{shervashidze2011wl}
& Shervashidze et al. 
& 2011 
& Efficient 2D graph kernel 
& Ignores 3D structural information \\ \hline

Pharmacophore Kernel \cite{mahe2006pharmacophore}
& Mahé et al. 
& 2006 
& Captures spatial geometry 
& Computationally expensive; incomplete torsion modeling \\ \hline

Graph Convolutional Networks \cite{kipf2016gcn}
& Kipf \& Welling 
& 2016 
& Graph-based deep learning model 
& Requires large labeled datasets and training \\ \hline

Efficient 3D Kernels \cite{ankit2023efficient3d}
& Ankit et al. 
& 2023 
& Uses bond length, angle, torsion 
& High computational cost for large-scale screening \\ \hline


\end{tabular}
\end{frame}

% --- TABLE SLIDE 2  ---
\begin{frame}
\scriptsize
\setlength{\tabcolsep}{2pt}

\begin{tabular}{|p{2.8cm}|p{2.3cm}|p{0.9cm}|p{3.2cm}|p{3.8cm}|}
\hline
\textbf{Paper Title} & \textbf{Authors} & \textbf{Year} & \textbf{Focus} & \textbf{Key Insight} \\ \hline

Graph Kernel Framework \cite{gartner2003graphkernels}
& Gärtner et al. 
& 2003 
& R-convolution graph decomposition 
& Not optimized for molecular geometry or scalability \\ \hline

Graph Kernels in Drug Discovery \cite{palivela2017graphkernels}
& Palivela et al. 
& 2017 
& Comparative study of graph kernels 
& No unified efficient hybrid 2D–3D approach \\ \hline

Interpretable Molecular Property Prediction using Graph Kernels \cite{xiang2023interpretable}
& Yan Xiang et al. 
& 2023 
& Interpretable graph kernel models 
& Improves explainability of kernel-based predictions \\ \hline

Motif Entropy Graph Kernel \cite{motif2023kernel}
& (Pattern Recognition Journal) 
& 2023 
& Motif + entropy based graph kernel 
& Captures deeper structural patterns than WL kernel \\ \hline

Fingerprint + Graph Hybrid Neural Network \cite{zhang2024hybrid}
& Qingtian Zhang et al. 
& 2024 
& Combines fingerprints + graph features 
& Hybrid models outperform single-representation methods \\ \hline

Graph + Transformer Fusion Model \cite{prakash2023fusion}
& Sai Prakash et al. 
& 2023 
& GNN + Transformer integration 
& Captures both local and global molecular features \\ \hline

Efficient Graph Similarity using Topological Indices \cite{dehmer2025topological}
& Dehmer et al. 
& 2025 
& Topological index-based similarity 
& Reduces computational complexity vs kernel methods \\ \hline

\end{tabular}

\end{frame}

% SLIDE ---4. Research gap ----
\begin{frame}{4. Overall Research Gap}
\begin{itemize}
    \item Existing methods either rely on 2D representations (e.g., WL Kernel), which are computationally efficient but fail to capture detailed 3D structural information, or on 3D/multimodal models, which are accurate but computationally expensive and not scalable.

    \item Recent deep learning approaches (GNNs, Transformer-based, hybrid models) require large training datasets and high computational resources, limiting their applicability for efficient large-scale virtual screening.

    \item Many approaches focus on either topological features (2D) or geometric features (3D), with a lack of effective and efficient integration of both representations.

    \item Existing graph kernel methods and similarity techniques often rely on dense representations (e.g., adjacency matrices), leading to increased computational complexity.

    \item Therefore, there is a need for a scalable, efficient, and unified framework that balances:
    \begin{itemize}
        \item Fast \textbf{2D screening}
        \item Accurate \textbf{3D structural similarity}
        \item Reduced \textbf{computational complexity}
        \item Effective \textbf{candidate prioritization}
    \end{itemize}
\end{itemize}
\end{frame}

% --- SLIDE 5: EXISTING MECHANISMS ---
\section{Existing Mechanisms}
\begin{frame}{5. Existing Mechanisms}
\framesubtitle{Methods for Ligand Similarity and Screening}
\begin{itemize}
    \item \textbf{Computational Approaches:}
    \begin{itemize}
        \item \textit{Random Walk} - Measures similarity by counting matching random walks (sequences of connected nodes) between two graphs. Suffers from high computational cost and the tottering problem (revisiting the same nodes reduces meaningful information).

        \item \textit{WLK (Weisfeiler--Lehman Kernel)} - Iteratively relabels nodes based on neighbors to capture hierarchical graph structure efficiently. Fails to capture continuous attributes and 3D geometry; limited to discrete topology.

        \item \textit{Pharmacophore Kernel} - Compares molecules based on spatial arrangements of chemical features (e.g., donor/acceptor groups). Computationally expensive for large datasets due to 3D matching complexity.

        \item \textit{3D Graph Kernel (3DGHK)} - Captures bond length, bond angle, and torsion angle using 3-hop neighborhoods to compute molecular similarity more effectively.
    \end{itemize}
\end{itemize}
\end{frame}

% --- SLIDE 6: PROPOSED APPROACH ---
\begin{frame}{6. Proposed Approach}
    \framesubtitle{A Two-Stage Graph Kernel Pipeline}

    \vspace{0.3em}

    \begin{columns}[T]

        \begin{column}{0.48\textwidth}
            \textbf{Stage 1 \hspace{0.2em} 2D Filtering}
            \vspace{0.3em}
            \hrule
            \vspace{0.4em}
            \begin{enumerate}
                \setlength\itemsep{0.5em}
                \item Represent molecules as graphs using \textbf{adjacency lists} 
                with 2D structural attributes
                \item \textbf{2D Kernel} computes similarity scores against 
                known active and inactive molecules
                \item Candidates with low 2D similarity are \textbf{discarded early}
            \end{enumerate}
        \end{column}

        \hfill\vrule\hfill

        \begin{column}{0.48\textwidth}
            \textbf{Stage 2 \hspace{0.2em} 3D Evaluation}
            \vspace{0.3em}
            \hrule
            \vspace{0.4em}
            \begin{enumerate}
                \setlength\itemsep{0.5em}
                \item \textbf{Promising candidates} from Stage 1 proceed to 
                3D evaluation
                \item \textbf{3D Kernel} evaluates 3D attributes on promising candidates only
                \item \textbf{ML/DL model} makes final prediction: 
                \textit{bind or not bind?}
            \end{enumerate}
        \end{column}

    \end{columns}

    \vspace{0.8em}
    \hrule
    \vspace{0.4em}

    \centering
    \small
    \textbf{2D filtering first} $\longrightarrow$ \textbf{3D evaluation only for promising candidates} 
    $\longrightarrow$ \textbf{Final prediction}
\end{frame}

%--------SLIDE 7: DATASET DESCRIPTION%
\begin{frame}{7. Dataset Description}

\textbf{Source:} Standard ligand-based virtual screening benchmark datasets

\begin{itemize}
    \item \textbf{Target Proteins:}
    \begin{itemize}
        \item The dataset comprises multiple protein targets.
        \item Each pair of files (actives/inactives) corresponds to a specific protein target.
    \end{itemize}

    \item \textbf{Data Format:}
    \begin{itemize}
        \item Stored in \textbf{Structure Data File (.sdf)} format.
        \item Encodes detailed molecular information, including:
        \begin{itemize}
            \item Atom types and bond connectivity
            \item 2D structural representation
            \item 3D spatial coordinates of atoms
        \end{itemize}
    \end{itemize}

    \item \textbf{Molecular Classes:}
    \begin{itemize}
        \item \textbf{Active compounds:} Experimentally validated ligands exhibiting binding affinity.
        \item \textbf{Inactive compounds:} Molecules with no significant binding activity.
    \end{itemize}

    \item \textbf{Role in Proposed Pipeline:}
    \begin{itemize}
        \item Active molecules serve as a \textbf{reference set} for similarity-based learning.
        \item Promising candidates undergo \textbf{3D kernel-based analysis} for precise similarity evaluation.
    \end{itemize}
\end{itemize}

\end{frame}

% --- SLIDE 8: ORIGINALITY (Rubric: 4 Marks) ---

\section{Originality}
\begin{frame}{8. Originality of Work}
\framesubtitle{Rubric: Novel Idea vs.\ Extension}

\begin{block}{What exists}
\begin{itemize}
    \item 2D kernel methods (e.g.\ WLK) are fast but discard all geometric information
    \item 3D similarity methods are accurate but expensive when applied to every candidate
    \item Hybrid 2D+3D models fuse both simultaneously --- accuracy improves but cost remains high
\end{itemize}
\end{block}

\vspace{0.3em}

\begin{block}{What we contribute}
\begin{itemize}
    \item A \textbf{staged pipeline} --- 2D screening using WCS-GNCCP with Frank-Wolfe gates candidates before 3D evaluation
    \item A \textbf{rank-one factorization} (Lemma~1) that provably eliminates one $O(n^3)$ sub-step per gradient, making each Frank-Wolfe iteration faster by a factor proportional to the molecule size
    \item Our 2D kernel outputs a \textbf{bijective atom correspondence map} alongside the score --- this reduces the 3D stage to a single $O(n)$ Kabsch alignment
    \item No existing method treats 2D as a \textbf{search-space reduction gate} before 3D --- that staging decision is our contribution
\end{itemize}
\end{block}

\end{frame}

% --- SLIDE 9 : Technical Knowledge --- 

\begin{frame}{9. Technical Knowledge (Algorithms/Models)}
\framesubtitle{Custom Two-Stage Graph-Based Screening Algorithm}

\begin{columns}[T]

    \begin{column}{0.48\textwidth}
        \textbf{Stage 1 --- Custom 2D Kernel} \hfill \textcolor{amritaMaroon}{\small $O(T \cdot n^3)$/pair}
        \vspace{0.3em}
        \hrule
        \vspace{0.4em}
        \begin{itemize}
            \setlength\itemsep{0.5em}
            \item Strip 3D coordinates from SDF --- retain atom type and bond connectivity only
            \item Build \textbf{weighted adjacency matrix} per molecule; Frobenius-normalise for scale invariance (Prop.~4)
            \item \textbf{WCS-GNCCP + Frank-Wolfe} optimisation produces similarity score $s \in [0,1]$ \textbf{and} bijective atom map $\mathbf{X}^*$
            \item Rank-one factorization (Lemma~1) reduces $\mathbf{U}$ from $O(n^3)$ to $O(n^2)$
            \item Discard candidates where $s < 0.7$ --- threshold is tunable
        \end{itemize}
    \end{column}

    \hfill\vrule\hfill

    \begin{column}{0.48\textwidth}
        \textbf{Stage 2 --- 3D Geometric Check} \hfill \textcolor{amritaMaroon}{\small $O(n)$}
        \vspace{0.3em}
        \hrule
        \vspace{0.4em}
        \begin{itemize}
            \setlength\itemsep{0.5em}
            \item Atom map $\mathbf{X}^*$ from Stage 1 already known --- \textbf{no combinatorial search repeated}
            \item \textbf{Kabsch alignment:} SVD of a fixed $3\times 3$ cross-covariance matrix --- $O(1)$; coordinate extraction and RMSD over $L$ atoms --- $O(n)$
            \item Hybrid kernel: $k_{\text{hybrid}} = (1{-}\beta)\,k_{\text{2D}} + \beta\,e^{-\gamma_{\text{3D}}\,\text{RMSD}}$
            \item Output: \textbf{Match / No Match} per candidate
        \end{itemize}
    \end{column}

\end{columns}

\vspace{0.4em}
\hrule
\vspace{0.2em}
\centering\small
Overall per-pair cost dominated by Stage~1; Stage~2 adds only microseconds.

\end{frame}

% --- SLIDE 10: CONCEPTUAL FLOW DIAGRAM ---
\begin{frame}{10.Conceptual / Algorithmic Flow}
\centering
\includegraphics[width=0.9\textwidth]{KernelApproachArchitectureDiagram.png}
\end{frame}

% --- SLIDE 11: PAPER / APPROACH ---
\begin{frame}{11. GNCCP-Based Weighted Common Subgraph Matching}

\centering

\vspace{0.25cm}

\includegraphics[width=0.34\textwidth]{gnccp.png}

\vspace{0.1cm}

\footnotesize
\textbf{Objective:}
$H_0(X)=\left\|A_{\mathrm{sub}G}-A_{\mathrm{sub}H}\right\|_F^2$

\end{frame}

% --- SLIDE: THEORETICAL GUARANTEE (Relabeling Invariance) ---
\section{Theoretical Guarantee}
\begin{frame}{12. Theoretical Guarantee}
\framesubtitle{Relabeling Invariance of the 2D Kernel Similarity Score}

\begin{block}{Proposition 1 (Relabeling Invariance)}
\small
Let $\mathbf{P}\in\{0,1\}^{M\times M}$, $\mathbf{Q}\in\{0,1\}^{N\times N}$ be permutation matrices encoding an arbitrary re-indexing of atoms in molecules $G,H$ (as stored in the SDF file), and let $\mathbf{A}_G' = \mathbf{P}\mathbf{A}_G\mathbf{P}^{T}$, $\mathbf{A}_H' = \mathbf{Q}\mathbf{A}_H\mathbf{Q}^{T}$. Then, for any atom correspondence $\mathbf{X}\in\mathcal{P}$,
\[
s\big(\mathbf{P}\mathbf{X}\mathbf{Q}^{T};\,\mathbf{A}_G',\mathbf{A}_H'\big) \;=\; s\big(\mathbf{X};\,\mathbf{A}_G,\mathbf{A}_H\big).
\]
\end{block}

\vspace{0.4em}

\begin{exampleblock}{Why it matters}
\small
\begin{itemize}
\setlength\itemsep{0.3em}
\item Atom order in an SDF file is arbitrary --- it reflects file-writing convention, not molecular identity
\item Guarantees our \textbf{2D kernel score} and \textbf{bijective atom map} $M:i\mapsto j$ are well-defined regardless of atom indexing
\item Makes Stage 1 filtering deterministic and reproducible across differently-ordered input files
\end{itemize}
\end{exampleblock}

\end{frame}
% ============================================================
% SLIDE 13: EXTENDED VALIDATION
% ============================================================

\begin{frame}{13. Extended Validation Across Application Domains}
\framesubtitle{Testing Proposition 1 Beyond the Synthetic Graph Example}

\begin{columns}[T]

% ------------------------------------------------------------
% DROSOPHILA
% ------------------------------------------------------------

\begin{column}{0.31\textwidth}

\textbf{Drosophila Brain Connectome}
\vspace{0.2em}
\hrule
\vspace{0.3em}

\small
\begin{itemize}
\setlength\itemsep{0.35em}

\item Reference connectome:
24 identified neurons across 8 neuropils

\item Sample connectome:
30 neurons with 20 corresponding neurons + 10 unmatched extras

\item Graph construction based on
\textbf{soma-to-soma distances}

\item Fly B subjected to
rigid rotation + translation

\item Partial matching:
$M=24,\;N=30,\;L=20$

\item \textbf{Result:}
16/20 correctly matched
(80\%) in the rectified experiment

\item Rotation/translation does not alter
the underlying distance structure

\end{itemize}

\end{column}


% ------------------------------------------------------------
% MOF
% ------------------------------------------------------------

\begin{column}{0.31\textwidth}

\textbf{Metal--Organic Framework Matching}
\vspace{0.2em}
\hrule
\vspace{0.3em}

\small
\begin{itemize}
\setlength\itemsep{0.35em}

\item Reference motif:
metal centers + organic linkers

\item Reference size:
$M=18$ atoms

\item Sample:
2 missing atoms + 6 decoy atoms

\item $65^\circ$ 3D rotation + translation

\item Matching:
$N=22,\;L=16$

\item Unary term $\mathbf{C}$:
metal $\leftrightarrow$ linker mismatch penalty

\item $\alpha=0.7$

\item \textbf{Result:}
16/16 atoms correctly matched
(100\%)

\item No metal--linker mismatches

\end{itemize}

\end{column}


% ------------------------------------------------------------
% SOCIAL NETWORK
% ------------------------------------------------------------

\begin{column}{0.31\textwidth}

\textbf{Social Network De-anonymization}
\vspace{0.2em}
\hrule
\vspace{0.3em}

\small
\begin{itemize}
\setlength\itemsep{0.35em}

\item Reference network:
22 users with scale-free structure

\item Sample:
4 users missing + 8 decoy accounts

\item Noisy / dropped edges

\item All user IDs re-numbered

\item Matching:
$N=30,\;L=18$

\item Unary term $\mathbf{C}$:
interest-category mismatch penalty

\item $\alpha=0.8$

\item \textbf{Result:}
10/18 users correctly identified
(55.6\%)

\item Sparse, low-degree nodes remain
intrinsically ambiguous

\end{itemize}

\end{column}

\end{columns}

\vspace{0.5em}

\centering
\small
\textit{The experiments test whether the matching objective is preserved
under changes that do not alter the underlying graph correspondence.}

\end{frame}


% ============================================================
% SLIDE 14: VALIDATION SUMMARY TABLE
% ============================================================

\begin{frame}{14. Empirical Validation of Proposition 1}
\framesubtitle{Consistency of Graph Matching Under Relabeling and Geometric Transformations}

\vspace{0.3em}

\begin{itemize}
\setlength\itemsep{0.35em}

\item Each domain was evaluated using the GNCCP-based WCS matching framework.

\item Graphs were subjected to vertex relabeling or
rigid geometric transformations while preserving the underlying structure.

\item Matching results were compared after translating predictions
back to the original labeling.

\item \textbf{Proposition 1 concerns the invariance of the WCS objective
under vertex relabeling; geometric invariance follows when the graph
weights themselves are distance-preserving.}

\end{itemize}

\vspace{0.5em}

\centering

\renewcommand{\arraystretch}{1.25}

\resizebox{\textwidth}{!}{%
\begin{tabular}{lcccc}
\toprule
\textbf{Domain}
&
\textbf{$(M,N,L)$}
&
\textbf{Perturbation}
&
\textbf{Accuracy}
&
\textbf{Same after undo?}
\\
\midrule

Paper Fig.~1 example
&
$(5,6,4)$
&
Vertex relabeling
&
Illustrative
&
Yes
\\

Drosophila connectome
&
$(24,30,20)$
&
Rotation + translation
&
80\% (16/20)
&
Yes
\\

MOF fragment
&
$(18,22,16)$
&
$65^\circ$ rotation + translation
&
100\% (16/16)
&
Yes
\\

Social network
&
$(22,30,18)$
&
Full ID relabeling
&
55.6\% (10/18)
&
Yes
\\

\bottomrule
\end{tabular}
}

\vspace{0.7em}

\begin{block}{Key Observation}
\small
Matching accuracy varies with graph ambiguity and noise,
but the \textbf{graph matching objective remains invariant to vertex
labeling}. Thus, changing vertex names does not change the underlying
WCS matching problem.
\end{block}

\end{frame}
% ============================================================
% SLIDE 15: MATH OPTIMIZATION 1
% ============================================================
\section{Mathematical \& Algorithmic Optimizations}
\begin{frame}{15. Mathematical Optimization: Complexity Reduction}
\framesubtitle{Rank-One Factorization and Honest Complexity Analysis}

\begin{columns}[T]
    \begin{column}{0.48\textwidth}
        \textbf{Lemma 1: $\mathbf{U} = \mathbf{r}\mathbf{r}^T$}
        \vspace{0.3em}
        \hrule
        \vspace{0.4em}
        \small
        \begin{itemize}
            \setlength\itemsep{0.4em}
            \item The co-matching indicator $\mathbf{U} = \mathbf{X}\mathbf{1}_{N\times N}\mathbf{X}^T$ na\"ively costs $O(MN^2{+}M^2N)$
            \item Factor: $\mathbf{1}_{N\times N} = \mathbf{1}_N\mathbf{1}_N^T$, so $\mathbf{U} = (X\mathbf{1}_N)(X\mathbf{1}_N)^T = \mathbf{r}\mathbf{r}^T$
            \item Cost: $O(MN + M^2)$ --- \textcolor{green!60!black}{\textbf{one cubic sub-step eliminated}}
        \end{itemize}
    \end{column}

    \hfill\vrule\hfill

    \begin{column}{0.48\textwidth}
        \textbf{Irreducible $O(N^3)$ Operations}
        \vspace{0.3em}
        \hrule
        \vspace{0.4em}
        \small
        \begin{itemize}
            \setlength\itemsep{0.4em}
            \item $\mathbf{X}\mathbf{A}_H\mathbf{X}^T$: bilinear form with dense $\mathbf{X}$ --- cannot be factored below $O(M^2N)$
            \item Hungarian assignment: inherently $O(\min(M,N)^3)$
            \item \textbf{Per-iteration:} $O(N^3)$ overall
        \end{itemize}
    \end{column}
\end{columns}

\vspace{0.5em}

\begin{block}{Measured Speed vs.\ Other Kernels (drug-like molecules, $N{\approx}30$)}
\centering\small
\renewcommand{\arraystretch}{1.15}
\begin{tabular}{lcr}
\toprule
\textbf{Kernel} & \textbf{Complexity} & \textbf{ms/pair} \\
\midrule
Weisfeiler--Lehman & $O(hN)$ & 0.1 -- 1 \\
Shortest Path & $O(N^2)$ & 1 -- 5 \\
\textbf{Our WCS-GNCCP (4,4)} & $\mathbf{T\!\cdot\!O(N^3)}$ & \textbf{$\sim$6} \\
Random Walk & $O(N^3)$ & 10 -- 30 \\
Optimal Assignment & $O(N^3)$ & 10 -- 50 \\
\bottomrule
\end{tabular}
\end{block}
\end{frame}

% ============================================================
% SLIDE 16: MATH OPTIMIZATION 2
% ============================================================
\begin{frame}{16. Algorithmic Optimization: Convergence \& Speed}
\framesubtitle{Frank-Wolfe Duality Gap and Vectorization}

\begin{columns}[T]
    \begin{column}{0.48\textwidth}
        \textbf{1. Frank-Wolfe Duality Gap}
        \hrule
        \vspace{0.4em}
        \small
        \begin{itemize}
            \item \textit{Before:} Termination relied on a simple Frobenius norm difference between iterations, which lacked mathematical guarantees.
            \item \textit{After:} We implemented the strict \textbf{Frank-Wolfe Duality Gap} condition ($g_t = \langle \nabla J, \mathbf{X} - \mathbf{Y} \rangle < \varepsilon$). This mathematically guarantees that the algorithm only stops when it reaches the true optimal topological match.
        \end{itemize}
    \end{column}

    \hfill\vrule\hfill

    \begin{column}{0.48\textwidth}
        \textbf{2. Vectorized Precomputation}
        \hrule
        \vspace{0.4em}
        \small
        \begin{itemize}
            \item \textit{Before:} Pairwise feature extractions caused massive redundant processing ($O(N^2)$ extractions).
            \item \textit{After:} Molecular graphs and adjacency matrices are strictly precomputed and cached.
            \item Feature comparisons were completely vectorized in NumPy, eliminating hundreds of nested loops per matching step.
        \end{itemize}
    \end{column}
\end{columns}
\end{frame}

% ============================================================
% SLIDE 17: BASELINE ACCURACY & THRESHOLDING
% ============================================================
\begin{frame}{17. 2D Kernel Baseline \& Thresholding Strategy}
\framesubtitle{Measured Performance and Speed Optimisation}

\begin{itemize}
    \item \textbf{Empirical Results:} The optimized 2D WCS-GNCCP kernel was benchmarked across molecular datasets (Aromatase, LBD, MMP, p53) containing known active and inactive compounds.
    \item \textbf{Performance:} 1-NN LOOCV classification accuracy of \textbf{80\%} on Aromatase (5+5 pilot); full benchmark range: \textbf{60--90\%} depending on target complexity.
    \item \textbf{Speed:} ${\sim}6$\,ms/pair with \texttt{removeHs=True} (heavy atoms only, $N{\approx}30$), competitive with Random Walk kernels.
    \item \textbf{Critical Speedup --- Hydrogen Removal:}
    \begin{itemize}
        \item With explicit H: $N{\approx}80$, ${\sim}27$\,ms/pair
        \item Without H: $N{\approx}33$, ${\sim}6$\,ms/pair --- \textbf{4--5$\times$ faster}
        \item Hydrogens are degree-1 leaf nodes that contribute zero topological information but cost $\sim(N_H/N_{\text{heavy}})^3 \approx 15\times$ compute
    \end{itemize}
    \item \textbf{Thresholding Strategy:} Similarity threshold $s > 0.7$ aggressively filters topologically distinct molecules, reducing the 3D candidate pool.
\end{itemize}
\end{frame}

% --- SLIDE 18: SUMMARY OF WORK COMPLETED ---
\begin{frame}{18. Summary of Work Completed}
\small
\begin{itemize}

\item Performed literature survey on graph kernels, virtual screening, and 3D molecular similarity methods

\item Defined the problem statement and proposed a two-stage 2D-to-3D filtering pipeline

\item Collected the benchmark dataset (active and inactive molecules in 3D SDF format) from Liu et al.\ 2023

\item Derived and proved the mathematical framework: Theorem~1 (structural gradient), Lemma~1 (rank-one factorization), Propositions~1--5 (GNCCP, kernel, normalization, projection)

\item \textbf{Implemented the full WCS-GNCCP engine} in Python/NumPy with optional CuPy GPU support --- all 18 lines of Algorithm~1 map to code

\item \textbf{Verified correctness:} smoke test (Benzene--Phenol, score 0.9955), full pipeline test (10 molecules, 1-NN 80\%, SVM functional)

\item \textbf{Benchmarked speed:} ${\sim}6$\,ms/pair at $(4,4)$ iterations with hydrogen removal --- competitive with Random Walk and Optimal Assignment kernels

\item Built Kaggle and Colab notebook generators for reproducible benchmarks

\end{itemize}
\end{frame}

% --- SLIDE 19: PLANS FOR UPCOMING SEMESTERS ---
\begin{frame}{19. Plans for Upcoming Semester}
\small
\begin{itemize}

\item[\checkmark] \textit{(Complete)} Develop and verify the 2D WCS-GNCCP similarity module with mathematical proofs and benchmarks

\item Implement the \textbf{3D alignment module}: use the bijective atom map $\mathbf{X}^*$ from Stage~1 as anchor points for Kabsch alignment ($O(n)$ per pair)

\item Compute \textbf{RMSD-based 3D kernel}: $k_{\text{3D}} = \exp(-\gamma_{\text{3D}} \cdot \text{RMSD})$ for the filtered candidate pairs

\item Design and evaluate the \textbf{hybrid kernel}: $k_{\text{hybrid}} = (1{-}\beta)\,k_{\text{2D}} + \beta\,k_{\text{3D}}$ with cross-validated $\beta$

\item Run \textbf{full-scale benchmarks} across all available protein targets (Aromatase, LBD, MMP, p53, and others) with SVM and 1-NN classifiers

\item Perform \textbf{ablation study}: 2D-only vs.\ 3D-only vs.\ hybrid, measuring accuracy and throughput

\item Compare results with existing methods (WL, Random Walk, 3DGHK) and document findings

\end{itemize}
\end{frame}

% ============================================================
% SLIDE 20: NEXT STEPS
% ============================================================
\begin{frame}{20. Next Steps: 2D to 3D Spatial Mapping}
\framesubtitle{Kabsch Alignment via the Bijective Atom Map}

\begin{block}{Bridging Topology and Geometry}
The 2D kernel outputs a bijective atom-to-atom matching matrix $\mathbf{X}^*$. This map serves as rigid anchor points for $O(n)$ 3D alignment --- no combinatorial search is repeated.
\end{block}

\vspace{0.3em}

\begin{columns}[T]
    \begin{column}{0.52\textwidth}
        \textbf{3D Alignment Pipeline}
        \vspace{0.2em}\hrule\vspace{0.3em}
        \small
        \begin{enumerate}
            \setlength\itemsep{0.3em}
            \item Extract matched 3D coordinates from $\mathbf{X}^*$ --- $O(L)$
            \item Compute centroids and center both point sets --- $O(L)$
            \item Cross-covariance $\mathbf{H} = \mathbf{P}^T\mathbf{Q}$ --- always $3{\times}3$ --- $O(L)$
            \item SVD of $\mathbf{H}$ --- $O(1)$ (fixed $3{\times}3$ matrix)
            \item Optimal rotation $\mathbf{R} = \mathbf{V}\mathbf{U}^T$ and RMSD --- $O(L)$
        \end{enumerate}
        \vspace{0.3em}
        \textbf{Total: $O(L) = O(n)$} per pair
    \end{column}

    \hfill\vrule\hfill

    \begin{column}{0.43\textwidth}
        \textbf{Hybrid Kernel}
        \vspace{0.2em}\hrule\vspace{0.3em}
        \small
        \[
            k_{\text{3D}} = e^{-\gamma_{\text{3D}} \cdot \text{RMSD}}
        \]
        \[
            k_{\text{hybrid}} = (1{-}\beta)\,k_{\text{2D}} + \beta\,k_{\text{3D}}
        \]
        \begin{itemize}
            \setlength\itemsep{0.3em}
            \item $\beta$ cross-validated per target
            \item 3D stage eliminates ``2D decoys'' (topological matches that clash in 3D space)
        \end{itemize}
    \end{column}
\end{columns}

\end{frame}

% --- REFERENCES ---
\section{References}
\begin{frame}[allowframebreaks]{References}
    \bibliographystyle{IEEEtran}
    \bibliography{ref2} 
\end{frame}

% --- THANK YOU ---
\begin{frame}
    \centering    
    {\Huge \textbf{\textcolor{amritaMaroon}{Thank You}}}
\end{frame}

\end{document}